% Reference model: https://oi-wiki.org/graph/mod-shortest-path/
% Independent implementation and validation: docs/validation.md.
\subsection{同余最短路}
\label{algo:mod-shortest-path}
\AlgorithmInfo
    {求正整数的非负整数组合能表示哪些数，并统计区间内不同可表示整数的个数。}
    {设步长数为 $k$、最小步长为 $p$；预处理 $O(kp\log(kp+2))$，空间 $O(kp+k)$；判定 $O(1)$，计数 $O(p)$。}
    {步长非空且均为正整数，小于 $4\times10^{18}$；$p<2^{31}$ 且状态规模可承受；余数从 0 开始，查询上界小于 $4\times10^{18}$。}

\subsubsection*{原理}
给定步长 $a_1,\ldots,a_k$，每种可以使用任意次，考虑
\[
    x=\sum_{i=1}^k c_i a_i,\qquad c_i\in\mathbb Z_{\ge0}.
\]
直接按数值做完全背包，状态数取决于查询上界。同余最短路取 $p=\min a_i$，只保留 $0,1,\ldots,p-1$ 共 $p$ 个余数状态；适合查询上界很大、最小步长较小的情况。

对每个余数 $r$ 和步长 $a_i$，建立转移
\[
    r\longrightarrow(r+a_i)\bmod p,\qquad w=a_i.
\]
从余数 0 出发，初始距离为 0。一条路径的边权和就是一种可表示整数，终点正是它的余数，所以最短距离 $dis_r$ 表示模 $p$ 余 $r$ 的最小可表示值。边权全为正，可以用 Dijkstra；模板按需枚举转移，不显式保存 $kp$ 条边。$p$ 的倍数只产生正权自环，不改善最短路，直接跳过。

因为 $p$ 本身也是可选步长，从 $dis_r$ 出发不断加 $p$，便可得到该余数类内所有更大的数。因此
\[
    x\text{ 可表示}\quad\Longleftrightarrow\quad
    dis_{x\bmod p}\le x\qquad(x\ge0).
\]
不同余数类互不重叠，区间 $[0,X]$ 内不同可表示整数的个数为
\[
    C(X)=\sum_{\substack{0\le r<p\\dis_r\le X}}
    \left(\left\lfloor\frac{X-dis_r}{p}\right\rfloor+1\right),
    \qquad C(X)=0\ (X<0).
\]
区间 $[L,R]$ 的答案为 $C(R)-C(L-1)$。这里统计的是不同整数，而非组合方案数；0 对应所有步长均使用 0 次，也计入可表示整数。

\subsubsection*{使用教程}
\InlineCode{build(a)} 传入普通的步长列表，一次完成初始化和最短路计算；再次调用会重建。\InlineCode{dis[r]} 保存余数 $r$ 的最小可表示值，\InlineCode{can(x)} 判断 $x$ 能否表示，\InlineCode{count(x)} 统计 $[0,x]$，\InlineCode{count(l,r)} 统计闭区间 $[l,r]$。查询必须在构建后进行，负数不计入答案。
\begin{lstlisting}[style=mtycpp,numbers=none]
ModShortestPath g;
g.build({6,9,20});
bool ok=g.can(29);
long long cnt=g.count(1,100);
long long low=g.dis[5];
\end{lstlisting}

\subsubsection*{需要注意}
若初始值为 $b$，后续只能增加这些步长，统计 $[L,R]$ 时查询 \InlineCode{count(L-b,R-b)}；平移需保证整数类型不溢出。初始值只平移答案，不改变步长建图。

若所有步长的最大公约数 $g>1$，部分余数不可达，也存在无限多个不可表示的正整数；不必强行假设所有余数都有最短路。若 $g=1$、$p>1$ 且全部距离均未被截断，最大不可表示整数为 $\max_r dis_r-p$；$p=1$ 时所有非负整数都可表示。

模板将距离不小于 \InlineCode{INF} 的状态截断；在约定查询范围内，这些状态不会贡献答案。\InlineCode{dis[r]==INF} 表示不可达或最小值超出该范围。\InlineCode{count(0,x)} 包含 0，只统计正整数时从 1 开始。

\needspace{7\baselineskip}
\subsubsection*{模板代码}
\CodeFile{code/graphs/mod_shortest_path.cpp}
